\documentclass[aspectratio=169]{beamer}

% =========================================================
% Кодировка и язык
% =========================================================
\usepackage[utf8]{inputenc}
\usepackage[T2A]{fontenc}
\usepackage[russian]{babel}

% =========================================================
% Математика
% =========================================================
\usepackage{amsmath,amssymb}
\usepackage{bm}

% =========================================================
% Графика и таблицы
% =========================================================
\usepackage{graphicx}
\usepackage{booktabs}
\usepackage{array}

% =========================================================
% TikZ (ER-диаграммы и схемы)
% =========================================================
\usepackage{tikz}
\usetikzlibrary{shapes.geometric, arrows.meta, positioning, calc}

% =========================================================
% Тема
% =========================================================
\usetheme{default}
\usecolortheme{default}

% =========================================================
% Цвета СВФУ
% =========================================================
\definecolor{svfuDark}{RGB}{10,40,90}
\definecolor{svfuAccent}{RGB}{30,144,255}
\definecolor{svfuLight}{RGB}{230,240,255}

% =========================================================
% Стили TikZ (ER-нотация Чена, схемы, уровни, циклы)
% =========================================================
\tikzset{
  entity/.style={
    rectangle, draw=svfuDark, very thick, fill=svfuLight,
    minimum width=2.4cm, minimum height=1cm,
    font=\bfseries\small, align=center, rounded corners=1pt
  },
  relationship/.style={
    diamond, draw=svfuDark, very thick, fill=svfuAccent!25,
    aspect=2.4, font=\scriptsize, align=center, inner sep=1pt
  },
  attribute/.style={
    ellipse, draw=svfuDark, thin, fill=white,
    minimum width=1.8cm, minimum height=0.55cm,
    font=\scriptsize, align=center
  },
  card/.style={font=\tiny, fill=svfuLight, inner sep=1pt},
  layer/.style={
    rectangle, draw=svfuDark, very thick, fill=svfuAccent!15,
    minimum width=6.2cm, minimum height=1.3cm,
    font=\bfseries\normalsize, align=center, rounded corners=2pt
  },
  phase/.style={
    rectangle, draw=svfuDark, thick, fill=svfuLight,
    rounded corners=4pt, minimum width=2.5cm, minimum height=1cm,
    font=\small\bfseries, align=center
  },
  catbox/.style={
    rectangle, draw=svfuDark, thick, fill=white,
    minimum width=2.6cm, minimum height=0.9cm,
    font=\footnotesize, align=center, rounded corners=1pt
  }
}

% =========================================================
% Основные цвета Beamer
% =========================================================
\setbeamercolor{normal text}{fg=black,bg=white}
\setbeamercolor{structure}{fg=svfuDark}
\setbeamercolor{alerted text}{fg=svfuDark}
\setbeamercolor{frametitle}{bg=svfuDark,fg=white}
\setbeamercolor{title}{fg=white}
\setbeamercolor{author}{fg=white}
\setbeamercolor{date}{fg=white}

\setbeamercolor{accent line}{bg=svfuAccent}

\setbeamercolor{footline left}
  {bg=svfuDark,fg=white}

\setbeamercolor{footline center}
  {bg=svfuDark!85,fg=white!80}

\setbeamercolor{footline right}
  {bg=svfuDark!70,fg=white}

\setbeamercolor{block title}
  {bg=svfuDark,fg=white}

\setbeamercolor{block body}
  {bg=svfuLight,fg=black}

% =========================================================
% Шрифты
% =========================================================
\setbeamerfont{frametitle}
  {size=\large,series=\bfseries}

\setbeamerfont{title}
  {size=\Large,series=\bfseries}

% =========================================================
% Убираем навигацию
% =========================================================
\setbeamertemplate{navigation symbols}{}
\setbeamertemplate{headline}{}

% =========================================================
% Заголовок слайда
% =========================================================
\setbeamertemplate{frametitle}{%
  \vspace*{-1pt}%
  \begin{beamercolorbox}[
    wd=\paperwidth,
    ht=3.2ex,
    dp=1.2ex,
    leftskip=10pt,
    sep=0pt
  ]{frametitle}
    \usebeamerfont{frametitle}
    \insertframetitle
  \end{beamercolorbox}%
}

% =========================================================
% Подвал
% =========================================================
\setbeamertemplate{footline}{%
  \begin{beamercolorbox}[
    wd=\paperwidth,
    ht=0.4pt,
    dp=0pt
  ]{accent line}
  \end{beamercolorbox}%

  \hbox to \paperwidth{%

    \begin{beamercolorbox}[
      wd=0.305\paperwidth,
      ht=3ex,
      dp=1.5ex,
      left,
      sep=0pt,
      leftskip=6ex
    ]{footline left}
      \textbf{СВФУ}
    \end{beamercolorbox}%

    \begin{beamercolorbox}[
      wd=0.295\paperwidth,
      ht=3ex,
      dp=1.5ex,
      center,
      sep=0pt
    ]{footline center}
    \end{beamercolorbox}%

    \begin{beamercolorbox}[
      wd=0.4\paperwidth,
      ht=3ex,
      dp=1.5ex,
      right,
      sep=0pt
    ]{footline right}
      \raggedleft
      \today
      \hspace{5ex}
      \insertframenumber\,/\,\inserttotalframenumber
      \hspace{6ex}
    \end{beamercolorbox}%
  }%
}

% =========================================================
% Настройки списков
% =========================================================
\setbeamertemplate{itemize item}
  {\color{svfuAccent}\small$\blacktriangleright$}

\setbeamertemplate{itemize subitem}
  {\color{svfuAccent}\scriptsize$\blacktriangleright$}

\setlength{\leftmargini}{1.5em}

% =========================================================
% Документ
% =========================================================
\begin{document}

% =========================================================
% 1. ТИТУЛ
% =========================================================
\begin{frame}[plain]
  \vspace*{1cm}
  \begin{center}
    
    \colorbox{svfuDark}{%
      \parbox[c][2.8cm][c]{0.86\paperwidth}{%
        \centering
        {\LARGE\color{white}
          Лекция 1. Введение в проектирование баз данных. \\ Данные и модели представления.
        }
      }%
    }

  \end{center}

  \vfill

  \centering
  {\normalsize Дисциплина: Методы и средства проектирования баз данных}\\[4pt]
 {\small
   Составитель: Монастырев Егорий Егорьевич, М-ВТ-25 
 }\\[4pt]

 

  {\small
    Якутск, 2026
  }

  \vfill

\end{frame}


% % =========================================================
% % О ДИСЦИПЛИНЕ
% % =========================================================
% \begin{frame}{О дисциплине}

% Курс посвящён методам организации и
% \textbf{проектирования современных баз данных}.

% \vspace{0.3cm}

% \begin{itemize}
%   \item природа данных и их жизненный цикл;
%   \item концептуальное и логическое проектирование;
%   \item реляционная модель данных и язык SQL;
%   \item ограничения целостности и нормализация;
%   \item физическое проектирование;
%   \item создание приложений, работающих с БД.
% \end{itemize}

% \end{frame}





% =========================================================
% МЕСТО КУРСА
% =========================================================
\begin{frame}{Введение в проектирование баз данных}

Проектирование баз данных — одна из фундаментальных
инженерных компетенций разработчика программного обеспечения:
веб-приложения, банковские и научные системы, корпоративное ПО.

\vspace{0.5cm}

\begin{alertblock}{Главная идея}
Разработчик работает не только с кодом,
но и с данными, которые этот код создаёт и использует.
Прежде чем спроектировать хранение, необходимо понять
природу самих данных.
\end{alertblock}

\end{frame}

% --- Контрольные вопросы ---
\begin{frame}{Ключевые вопросы}

\small
\begin{enumerate}
  \item Чем данные отличаются от информации и знаний?
  \item Как классифицируются данные по источнику, формату и структуре?
  \item Из каких этапов состоит жизненный цикл данных?
  \item Чем БД отличается от СУБД?
  \item Чем табличное представление отличается от документного?
  \item Чем JSON отличается от NoSQL?
  \item Что такое объектное хранилище?
  \item Что такое ER-диаграмма и из чего она состоит?
  \item Что такое архитектура ANSI/SPARC?
  \item Зачем нужна нормализация?
\end{enumerate}

\end{frame}


% % =========================================================
% % ЦЕЛЬ ЛЕКЦИИ
% % =========================================================
% \begin{frame}{Цель лекции}

% \begin{itemize}
%   \item Разграничить понятия «данные», «информация», «знания».
%   \item Понять принципы классификации и жизненный цикл данных.
%   \item Понять, что такое база данных и СУБД, зачем нужно проектирование.
%   \item Сравнить табличное, документное и файловое представление данных.
%   \item Освоить базовые понятия ER-моделирования (сущность, атрибут, связь).
%   \item Увидеть путь от предметной области до физической реализации в PostgreSQL.
% \end{itemize}

% \vspace{0.3cm}

% \begin{block}{Результат}
% После лекции студент должен видеть БД
% не как набор таблиц, а как результат
% последовательного инженерного процесса анализа данных.
% \end{block}

% \end{frame}


% =========================================================
% ЧАСТЬ I. ДАННЫЕ И ИНФОРМАЦИЯ
% =========================================================

% --- Данные: определение ---
\begin{frame}{Данные: определение}

В стандартах и терминологических словарях понятие данных
связывается с представлением информации в формализованной
форме, пригодной для передачи, интерпретации или обработки.

\vspace{0.3cm}

\begin{block}{Определение (по терминологии NIST)}
\textbf{Данные} — представление фактов, концепций
или инструкций в форме, пригодной для коммуникации,
интерпретации или обработки человеком либо
автоматическими средствами.
\end{block}

\vspace{0.3cm}

Примеры данных:

\begin{center}
\texttt{101} \quad
\texttt{Иванов И.И.} \quad
\texttt{23.7} \quad
\texttt{2026-09-02}
\end{center}

\end{frame}


% --- Семантика данных ---
\begin{frame}{Семантика данных}

Само значение не всегда содержит информацию о том,
что именно оно обозначает.

\begin{center}
\Huge
\texttt{101}
\end{center}

\begin{center}
может быть:
\end{center}

\begin{center}
идентификатор студента
\quad
номер группы
\quad
номер заказа
\quad
код объекта
\end{center}

\vspace{0.3cm}

\begin{alertblock}{Вывод}
Значение данных должно рассматриваться совместно
с контекстом и правилами его интерпретации.
\end{alertblock}

\end{frame}


% --- Данные, информация, знания ---
\begin{frame}{Данные, информация, знания}

Для учебных целей полезно разграничивать три уровня:

\vspace{0.3cm}

\begin{center}
\Large
Данные $\rightarrow$ Информация $\rightarrow$ Знания
\end{center}

\vspace{0.4cm}

\begin{itemize}
  \item \textbf{Данные} — формализованные представления
        фактов, наблюдений, измерений.
  \item \textbf{Информация} — данные, интерпретированные
        в определённом контексте и обладающие смыслом
        для получателя (согласуется с терминологией ISO и NIST).
  \item \textbf{Знания} — организованная совокупность понятий
        и связей между ними, позволяющая интерпретировать
        информацию и использовать её при решении задач.
\end{itemize}

\vspace{0.2cm}
\footnotesize
Термин «знание», в отличие от «данных» и «информации»,
имеет разные трактовки в литературе; здесь он используется
как более высокий уровень организации информации.

\end{frame}


\begin{frame}{Пример: от данных к знаниям}

\begin{columns}[T]
\column{0.32\textwidth}
\begin{block}{\centering Данные}
\begin{center}
\small
\texttt{Студент: 101}\\[0.1cm]
\texttt{Математика: 5}\\[0.1cm]
\texttt{Программирование: 5}\\[0.1cm]
\texttt{Физика: 3}
\end{center}
\end{block}

\column{0.32\textwidth}
\begin{block}{\centering Информация}
\small
Студент 101 имеет высокие оценки по математике и программированию и более низкую — по физике.
\end{block}

\column{0.32\textwidth}
\begin{block}{\centering Знание}
\small
Выраженная успеваемость студента в дисциплинах, связанных с математикой и программированием.
\end{block}
\end{columns}

\vspace{0.7cm}

\begin{center}
\large
\textbf{Данные}
\;$\longrightarrow$\;
\textit{Интерпретация}
\;$\longrightarrow$\;
\textbf{Информация}
\;$\longrightarrow$\;
\textit{Обобщение опыта}
\;$\longrightarrow$\;
\textbf{Знание}
\end{center}

\end{frame}


\begin{frame}{Почему важно различать данные и информацию}

Проектирование информационной системы предполагает  
работу не просто со значениями, а с их:

\vspace{0.4cm}

\begin{center}
\begin{tabular}{ccc}
\textbf{смыслом} & \textbf{структурой} & \textbf{происхождением} \\[0.3cm]
\textbf{взаимосвязями} & \textbf{ограничениями} & \textbf{временем существования}
\end{tabular}
\end{center}

\vspace{0.7cm}

\begin{alertblock}{Ключевой вопрос проектировщика}
\begin{center}
Не только \Large\textbf{что хранить?}\\[0.3cm]
но и \Large\textbf{что означает каждый элемент данных?}
\end{center}
\end{alertblock}

\end{frame}


% =========================================================
% ЧАСТЬ II. КЛАССИФИКАЦИЯ ДАННЫХ
% =========================================================

% --- Обзорная схема классификации (TikZ) ---
\begin{frame}{Классификация данных: обзор}

\begin{center}
\begin{tikzpicture}[node distance=0.9cm and 1.6cm]

\node[entity, fill=svfuDark, text=white] (root) {ДАННЫЕ};

\node[catbox, below=1.3cm of root, xshift=-4.2cm] (src) {по источнику};
\node[catbox, below=1.3cm of root] (fmt) {по формату};
\node[catbox, below=1.3cm of root, xshift=4.2cm] (str) {по структуре};

\draw[thick, svfuDark] (root) -- (src);
\draw[thick, svfuDark] (root) -- (fmt);
\draw[thick, svfuDark] (root) -- (str);

\node[attribute, below=0.5cm of src, xshift=-0.9cm, font=\small] (s1) {наблюдение};
\node[attribute, below=1.3cm of src, xshift=0.9cm, font=\small] (s2) {эксперимент};
\draw[thin] (src) -- (s1); \draw[thin] (src) -- (s2);

\node[attribute, below=0.5cm of fmt, xshift=-0.9cm, font=\small] (f1) {текст, числа};
\node[attribute, below=1.3cm of fmt, xshift=0.9cm, font=\small] (f2) {аудио, видео};
\draw[thin] (fmt) -- (f1); \draw[thin] (fmt) -- (f2);

\node[attribute, below=0.5cm of str, xshift=-0.9cm, font=\small] (t1) {структ.};
\node[attribute, below=1.3cm of str, xshift=0.9cm, font=\small] (t2) {неструкт.};
\draw[thin] (str) -- (t1); \draw[thin] (str) -- (t2);

\end{tikzpicture}
\end{center}

Дополнительно данные классифицируются
по характеру изменения во времени.

\end{frame}


% --- По источнику ---
\begin{frame}{Откуда берутся данные}

\small

\begin{table}

\centering

\begin{tabular}{>{\bfseries}p{0.3\textwidth}|p{0.55\textwidth}}

\toprule

Источник & Примеры \\

\midrule

Наблюдения & датчики, опросы, измерения \\

Эксперименты & результаты лабораторных экспериментов \\

Моделирование & результаты математических моделей \\

Результаты обработки & статистика, средние значения, прогнозы \\

Справочные данные & классификаторы, справочники \\

\bottomrule

\end{tabular}

\end{table}

\vspace{0.4cm}

\begin{center}

\large

Данные $\rightarrow$ Обработка $\rightarrow$ Новые данные

\end{center}

\end{frame}


% --- По формату ---

\begin{frame}{В каком виде представлены данные}

\begin{itemize}

    \item \textbf{Текст} — документы, сообщения, описания;
    
    \item \textbf{Числа} — измерения, оценки, статистика;
    
    \item \textbf{Изображения} — фотографии, схемы, рисунки;
    
    \item \textbf{Аудио и видео} — звукозаписи, фильмы, видеозаписи;
    
    \item \textbf{Файлы} — PDF, архивы, программы и другие данные.

\end{itemize}

\vspace{0.4cm}

\begin{alertblock}{Важно}

Формат показывает, \textbf{в каком виде записаны данные}, но не объясняет их смысл.

\end{alertblock}

\end{frame}


% --- Структурированность ---

\begin{frame}{Как организованы данные}

\small

\begin{table}

\centering

\begin{tabular}{>{\bfseries}p{0.28\textwidth}|p{0.30\textwidth}|p{0.30\textwidth}}

\toprule

Тип & Что это значит & Примеры \\

\midrule

Структурированные &
имеют заранее заданную структуру &
таблицы, базы данных \\

Полуструктурированные &
структура есть, но она более гибкая &
JSON, XML \\

Неструктурированные &
нет заранее заданной структуры &
текст, фото, видео \\

\bottomrule

\end{tabular}

\end{table}

\vspace{0.4cm}

\begin{alertblock}{Главное}

Чем сложнее структура данных, тем больше вариантов их хранения и обработки.

\end{alertblock}

\end{frame}


\begin{frame}{Как изменяются данные со временем}

\begin{columns}[T]
\column{0.32\textwidth}
\begin{block}{\centering Неизменяемые}
\small
После создания данные \textbf{не меняются}.

\vspace{0.25cm}
\textit{Примеры:}\\
архив,\\
результат эксперимента
\end{block}

\column{0.32\textwidth}
\begin{block}{\centering Растущие}
\small
Старые данные сохраняются,\\
добавляются \textbf{новые}.

\vspace{0.25cm}
\textit{Примеры:}\\
журнал событий,\\
история измерений
\end{block}

\column{0.32\textwidth}
\begin{block}{\centering Изменяемые}
\small
Существующие данные\\
могут \textbf{изменяться} или удаляться.

\vspace{0.25cm}
\textit{Примеры:}\\
профиль пользователя
\end{block}
\end{columns}

\vspace{0.6cm}

\begin{alertblock}{Почему это важно?}
Способ изменения данных напрямую влияет на их хранение, обновление и резервное копирование.
\end{alertblock}

\end{frame}


\begin{frame}{Большие данные (Big Data)}

\begin{block}{Определение}
Термин \textbf{Big Data} обозначает данные, 
\textbf{объём}, \textbf{разнообразие}, \textbf{скорость поступления}
или другие характеристики которых создают существенные
требования к традиционным методам их обработки и хранения.
\end{block}

\vspace{0.6cm}

\begin{center}
\large\textbf{Примеры источников}
\end{center}

\vspace{0.2cm}

\begin{columns}[T]

\column{0.32\textwidth}
\begin{itemize}
    \item транзакции
    \item журналы событий
\end{itemize}

\column{0.32\textwidth}
\begin{itemize}
    \item сенсорные данные
    \item социальные данные
\end{itemize}

\column{0.32\textwidth}
\begin{itemize}
    \item изображения
    \item видео
\end{itemize}

\end{columns}

\end{frame}


% =========================================================
% ЧАСТЬ III. ЖИЗНЕННЫЙ ЦИКЛ И ЗАЩИТА ДАННЫХ
% =========================================================

% --- Жизненный цикл (TikZ, цикл) ---
\begin{frame}{Жизненный цикл данных}

Данные существуют не только в момент их записи в базу.

\begin{center}
\begin{tikzpicture}[node distance=2.2cm]

\node[phase] (plan) at (90:2.6cm) {Планирование};
\node[phase] (collect) at (18:2.6cm) {Сбор};
\node[phase] (process) at (-54:2.6cm) {Обработка\\и анализ};
\node[phase] (preserve) at (-126:2.6cm) {Сохранение\\и публикация};
\node[phase] (reuse) at (162:2.6cm) {Повторное\\использование};

\draw[-{Latex[length=2.5mm]}, thick, svfuAccent] (plan) -- (collect);
\draw[-{Latex[length=2.5mm]}, thick, svfuAccent] (collect) -- (process);
\draw[-{Latex[length=2.5mm]}, thick, svfuAccent] (process) -- (preserve);
\draw[-{Latex[length=2.5mm]}, thick, svfuAccent] (preserve) -- (reuse);
\draw[-{Latex[length=2.5mm]}, thick, svfuAccent] (reuse) -- (plan);

\end{tikzpicture}
\end{center}

\footnotesize
\begin{center}
Согласно UK Data Service, жизненный цикл исследовательских
данных представляется этапами:
plan~$\rightarrow$~collect~$\rightarrow$~process and analyse~$\rightarrow$~preserve and share~$\rightarrow$~reuse.
\end{center}

\end{frame}


% --- Этапы кратко ---
\begin{frame}{Ключевые этапы жизненного цикла}

\small
\begin{table}
\centering
\begin{tabular}{>{\bfseries}p{0.28\textwidth}|p{0.6\textwidth}}
\toprule
Этап & Суть\\
\midrule
Планирование &
определяются форматы, место хранения, доступ,
требования к защите — \textit{до} создания данных\\
Сбор &
данные регистрируются, классифицируются,
снабжаются метаданными\\
Обработка и анализ &
проверка, очистка, преобразование, агрегация;
результат — производные данные\\
Документирование &
сохраняются сведения о происхождении, структуре,
методике получения, версиях\\
Сохранение и доступ &
резервные и архивные копии, контроль целостности,
управление правами доступа\\
Повторное использование &
хорошо документированные данные используются
в новых задачах\\
\bottomrule
\end{tabular}
\end{table}

\end{frame}


% --- Защита данных: CIA ---
\begin{frame}{Защита данных: конфиденциальность, целостность, доступность}

\begin{center}
\begin{tikzpicture}[node distance=1.6cm and 1.6cm]
\node[entity, fill=svfuDark, text=white] (root) {Защита данных};
\node[catbox, below=1.2cm of root, xshift=-3.6cm] (c) {Конфиденциальность};
\node[catbox, below=1.2cm of root] (i) {Целостность};
\node[catbox, below=1.2cm of root, xshift=3.6cm] (a) {Доступность};
\draw[thick, svfuDark] (root) -- (c);
\draw[thick, svfuDark] (root) -- (i);
\draw[thick, svfuDark] (root) -- (a);
\end{tikzpicture}
\end{center}

\begin{itemize}
  \item \textbf{Конфиденциальность} — доступ получают только
        уполномоченные субъекты;
  \item \textbf{Целостность} — данные не должны изменяться
        несанкционированным образом;
  \item \textbf{Доступность} — уполномоченные пользователи
        должны получать данные тогда, когда это необходимо.
\end{itemize}

\end{frame}


% =========================================================
% ЧАСТЬ IV. БАЗА ДАННЫХ И СУБД
% =========================================================

% --- Что такое БД ---
\begin{frame}{Что такое база данных}

\begin{block}{Определение (по терминологии NIST)}
\textbf{База данных} рассматривается как репозиторий
информации или данных; такое определение не ограничивает
базу данных только традиционной реляционной системой.
\end{block}

\vspace{0.4cm}

\begin{alertblock}{Основное положение}
База данных не тождественна таблице.
Таблица — лишь один из возможных способов
представления структурированных данных.
\end{alertblock}

\end{frame}


% --- БД в АИС ---
\begin{frame}{База данных как основа информационной системы}

Современные информационные системы работают
со структурированными данными: интернет-магазин,
банковская система, корпоративная CRM, социальная сеть.

\vspace{0.4cm}

\begin{center}
\Large
\textbf{Приложение}
$\longleftrightarrow$
\textbf{База данных}
\end{center}

\vspace{0.4cm}

\textbf{АИС} (автоматизированная информационная система) —
совокупность программно-технических средств и персонала,
обеспечивающая сбор, хранение, обработку и предоставление
информации. БД — важнейший, но не единственный компонент АИС.

\end{frame}


% --- Наивный подход ---
\begin{frame}{Наивный подход и его проблемы}

\begin{center}
\large
Создать таблицы
$\rightarrow$
Добавить столбцы
$\rightarrow$
Записать данные
\end{center}

\vspace{0.3cm}

На небольшом примере такой подход может работать,
но при росте системы появляются проблемы:

\begin{itemize}
  \item дублирование и противоречивые значения;
  \item аномалии вставки, обновления и удаления;
  \item сложность изменения структуры;
  \item падение производительности запросов.
\end{itemize}

\vspace{0.3cm}

\begin{alertblock}{Главная мысль}
Ошибки в структуре данных часто обнаруживаются
только после того, как система уже начала работать.
\end{alertblock}

\end{frame}


% --- Зачем нужно проектирование ---
\begin{frame}{Почему проектирование необходимо?}

\begin{center}
\Large
\textbf{Проектирование БД $\neq$ CREATE TABLE}
\end{center}

\vspace{0.5cm}

Проектирование отвечает на вопросы: какие данные хранить,
как они связаны, какие ограничения существуют,
какие операции будут выполняться и как система
будет расти и масштабироваться.

\vspace{0.4cm}

\begin{block}{Определение}
\textbf{Проектирование базы данных} — процесс определения
структуры, взаимосвязей, ограничений и способов хранения
данных в соответствии с требованиями информационной системы.
\end{block}

\end{frame}


% --- БД и СУБД ---
\begin{frame}{База данных и СУБД}

\begin{table}
\centering
\begin{tabular}{>{\bfseries}p{0.25\textwidth}|p{0.62\textwidth}}
\toprule
Понятие & Определение\\
\midrule
База данных &
структурированный (или иначе организованный) репозиторий данных\\
СУБД &
программное обеспечение для создания,
хранения и обработки БД\\
\bottomrule
\end{tabular}
\end{table}

\vspace{0.4cm}

Примеры СУБД:

\begin{center}
PostgreSQL \quad Microsoft Access \quad MySQL \quad Oracle Database
\end{center}

\end{frame}


% --- Функции СУБД ---
\begin{frame}{Зачем нужна СУБД и что она делает}

Просто хранить файлы с данными недостаточно.

\vspace{0.2cm}
\small
\begin{table}
\centering
\begin{tabular}{p{0.32\textwidth}|p{0.52\textwidth}}
\toprule
\textbf{Функция} & \textbf{Суть}\\
\midrule
Хранение & таблицы, индексы, страницы\\
Транзакции & атомарное выполнение операций\\
Журнализация & восстановление после сбоев\\
Целостность & проверка ограничений\\
Безопасность & управление доступом и правами\\
Оптимизация & эффективное выполнение запросов\\
\bottomrule
\end{tabular}
\end{table}

\end{frame}


% --- Метаданные и словарь данных ---
\begin{frame}{Метаданные и словарь данных}

\textbf{Метаданные} — данные, описывающие структуру,
свойства и правила использования других данных
(«данные о данных»): имена таблиц и столбцов,
типы данных, ключи, ограничения, права доступа.

\vspace{0.4cm}

\textbf{Словарь данных} (Data Dictionary) —
компонент СУБД, содержащий метаданные о структуре БД.

\vspace{0.2cm}
В PostgreSQL:

\begin{center}
\texttt{pg\_catalog}
\qquad
\texttt{information\_schema}
\end{center}

\end{frame}


% =========================================================
% ЧАСТЬ V. МОДЕЛИ ПРЕДСТАВЛЕНИЯ И ХРАНЕНИЯ ДАННЫХ
% =========================================================

% --- Три подхода: обзор (TikZ) ---
\begin{frame}{Способы представления и хранения данных}

\begin{center}
\begin{tikzpicture}[node distance=1.2cm and 1.4cm]
\node[entity, fill=svfuDark, text=white] (root) {Данные};
\node[catbox, below=1.3cm of root, xshift=-4.5cm] (t) {Табличное\\представление};
\node[catbox, below=1.3cm of root, xshift=-1.40cm] (d) {Документное\\представление};

\node[catbox, below=1.3cm of root, xshift=1.60cm] (f) {Файловое\\хранение};
\node[catbox, below=1.3cm of root, xshift=4.5cm]
    (v) {Векторное\\представление};

\draw[thick, svfuDark] (root) -- (v);

\node[attribute, below=0.6cm of v, font=\scriptsize]
    (v2) {Vector DB};

\draw[thin] (v) -- (v2);
\draw[thick, svfuDark] (root) -- (t);
\draw[thick, svfuDark] (root) -- (d);
\draw[thick, svfuDark] (root) -- (f);
\node[attribute, below=0.6cm of t, font=\scriptsize] (t2) {SQL};
\node[attribute, below=0.6cm of d, font=\scriptsize] (d2) {JSON};
\node[attribute, below=0.6cm of f, font=\scriptsize] (f2) {Object Storage};
\draw[thin] (t) -- (t2); \draw[thin] (d) -- (d2); \draw[thin] (f) -- (f2);

\end{tikzpicture}
\end{center}

Важно различать \textbf{модель представления данных}
и \textbf{физическое средство хранения} — это не одно и то же.

\end{frame}


% --- Реляционное представление ---
\begin{frame}{Реляционное представление данных}

В реляционной модели данные представлены
в виде \textbf{отношений}; практическим представлением
отношения является таблица.

\vspace{0.3cm}

\begin{center}
\begin{tabular}{|c|c|c|}
\hline
id & name & group\_id\\
\hline
1 & Иванов & 101\\
2 & Петров & 101\\
3 & Сидоров & 102\\
\hline
\end{tabular}
\end{center}

\vspace{0.3cm}

Табличная модель характеризуется фиксированной структурой,
атрибутами, идентификацией экземпляров, связями
и ограничениями целостности.

\end{frame}


% --- JSON ---
\begin{frame}[fragile]{JSON: определение и структура}

\begin{block}{Определение (RFC 8259)}
JSON — лёгкий текстовый, независимый от языка формат
обмена данными, предназначенный для переносимого
представления структурированных данных.
\end{block}

\vspace{0.3cm}
\small
\begin{verbatim}
{
  "id": 101,
  "name": "Иванов Иван",
  "group": { "number": "ПМИ-25", "course": 2 },
  "subjects": ["Математика", "Программирование"]
}
\end{verbatim}
\normalsize

JSON поддерживает объекты, массивы, строки, числа,
логические значения и \texttt{null}.

\end{frame}


% --- NoSQL vs JSON ---
\begin{frame}{NoSQL и JSON: важное различие}

\textbf{NoSQL} — широкий класс систем управления данными,
не ограниченных классической реляционной моделью.
В курсе рассматривается прежде всего документная модель
как один из вариантов нереляционного представления.

\vspace{0.5cm}

\begin{alertblock}{Принципиально важно}
JSON и NoSQL — не синонимы.\\[2pt]
JSON — \textbf{формат} представления данных.\\
NoSQL — более широкий \textbf{класс подходов}
к организации и управлению данными.
\end{alertblock}

\end{frame}


% --- Документное представление ---
\begin{frame}{Документное представление}

При документном подходе совокупность связанных атрибутов
объекта представляется единым документом с вложенной
структурой.

\vspace{0.3cm}

\begin{itemize}
  \item допускает различия в структуре отдельных документов;
  \item естественно представляет иерархические данные;
  \item приближает представление данных к структуре
        объектов приложения.
\end{itemize}

\vspace{0.3cm}

\begin{alertblock}{Оборотная сторона}
Гибкость структуры переносит часть ответственности
за согласованность данных на уровень проектирования
и программного обеспечения.
\end{alertblock}

\end{frame}

% --- Векторное представление ---

\begin{frame}{Векторное представление данных}

Некоторые объекты удобно описывать не набором отдельных полей, а \textbf{вектором признаков} в многомерном пространстве.

\[\Large
\mathbf{x} =
(x_1,x_2,\ldots,x_n)
\]
\normalsize
В задачах машинного обучения и поиска вектор может представлять:

\begin{itemize}

    \item текстовый документ;

    \item изображение;

    \item аудиозапись;

    \item профиль пользователя;

    \item результат работы нейронной сети.

\end{itemize}

\end{frame}


\begin{frame}{Сравнение способов представления данных}

\small

\begin{table}

\centering

\begin{tabular}{p{0.23\textwidth}|p{0.21\textwidth}|p{0.21\textwidth}|p{0.21\textwidth}}

\toprule

Характеристика &
Табличное &
Документное &
Векторное \\

\midrule

Основная единица &
строка /
отношение &
документ &
вектор \\

Структура &
явная схема &
гибкая /
иерархическая &
фиксированная
размерность \\

Основная операция &
запрос по атрибутам &
доступ к документу &
поиск по близости \\

Типичный пример &
PostgreSQL &
MongoDB &
Vector DB \\

\bottomrule

\end{tabular}

\end{table}

\vspace{0.3cm}

\begin{alertblock}{Важно}

Это не взаимоисключающие технологии. Одна информационная система может одновременно использовать таблицы, документы, файлы и векторные представления.

\end{alertblock}

\end{frame}


% --- Файловое и объектное хранение ---
\begin{frame}{Файловое и объектное хранение}

Для больших бинарных объектов (фото, видео, аудио, PDF,
архивы) используются файловые или объектные хранилища.

\vspace{0.3cm}

\begin{block}{Определение (по документации Amazon S3)}
Объект объектного хранилища включает как минимум
\textbf{содержимое} и \textbf{ключ}, по которому он
идентифицируется; также могут использоваться
\textbf{метаданные}. Объекты размещаются в контейнерах
(bucket).
\end{block}

\vspace{0.3cm}

\begin{center}
Bucket $\rightarrow$ Object (Key + Data + Metadata)
\end{center}

\end{frame}


% --- Один объект — несколько структур ---
\begin{frame}{Один объект — несколько структур хранения}

\begin{center}
\begin{tikzpicture}[node distance=1cm and 1.6cm]
\node[entity] (s) {Студент};
\node[catbox, below=1.1cm of s, xshift=-3.6cm] (a) {id, ФИО, группа};
\node[catbox, below=1.1cm of s] (b) {настройки};
\node[catbox, below=1.1cm of s, xshift=3.6cm] (c) {фотография};
\draw[thick] (s) -- (a); \draw[thick] (s) -- (b); \draw[thick] (s) -- (c);
\node[attribute, below=0.6cm of a, font=\scriptsize] (a2) {таблица};
\node[attribute, below=0.6cm of b, font=\scriptsize] (b2) {JSON};
\node[attribute, below=0.6cm of c, font=\scriptsize] (c2) {файл / объект};
\draw[thin] (a) -- (a2); \draw[thin] (b) -- (b2); \draw[thin] (c) -- (c2);
\end{tikzpicture}
\end{center}

\begin{alertblock}{Вывод}
Один объект предметной области может быть представлен несколькими структурами хранения одновременно.
\end{alertblock}

\end{frame}


% --- Критерии выбора хранения ---
\begin{frame}{Критерии выбора способа хранения}

\begin{itemize}
  \item структура данных: фиксированная или изменяемая;
  \item вложенность: плоская или иерархическая;
  \item объём: KB, MB, GB, TB;
  \item частота изменения: редкая или частая;
  \item характер доступа: по идентификатору, по атрибутам,
        по связям, последовательное чтение;
  \item требования к целостности: строгие ограничения
        или более гибкая структура.
\end{itemize}

\vspace{0.3cm}

\begin{block}{Вывод}
Место и способ хранения выбираются исходя
из свойств данных, а не наоборот.
\end{block}

\end{frame}


% =========================================================
% ЧАСТЬ VI. ПРЕДМЕТНАЯ ОБЛАСТЬ И ER-МОДЕЛИРОВАНИЕ
% =========================================================

% --- Предметная область ---
\begin{frame}{Предметная область}

\textbf{Предметная область} — часть реального мира
(фрагмент, рассматриваемый в рамках конкретной задачи
моделирования), информацию о которой необходимо
хранить и обрабатывать.

\vspace{0.4cm}

\begin{center}
\Large
университет \quad магазин \quad библиотека \quad больница
\end{center}

\vspace{0.4cm}

Первый шаг проектирования — \textbf{не SQL},
а анализ предметной области: какие объекты существуют,
какими характеристиками обладают, какие действия
выполняются, как объекты связаны.

\end{frame}


% --- Сущность / Атрибут / Связь ---
\begin{frame}{Сущность, атрибут, связь}

\begin{itemize}
  \item \textbf{Сущность} — значимый объект предметной
        области, сведения о котором должны быть
        представлены в информационной системе.
  \item \textbf{Атрибут} — характеристика сущности.
  \item \textbf{Связь} — отношение между сущностями
        предметной области.
\end{itemize}

\vspace{0.4cm}

\begin{center}
\Large
Студент
$\xrightarrow{\text{обучается в}}$
Группа
\end{center}

\begin{center}
\Large
Студент
$\xrightarrow{\text{получает}}$
Оценка
\end{center}

\end{frame}


% --- Концептуальная модель ---
\begin{frame}{Концептуальная модель}

Концептуальная модель описывает предметную область
\textbf{независимо от конкретной СУБД}.

\vspace{0.4cm}

На этом этапе неважно:

\begin{itemize}
  \item PostgreSQL или Access;
  \item какой тип данных используется;
  \item каким SQL-запросом это будет реализовано.
\end{itemize}

\vspace{0.4cm}

\begin{alertblock}{Главное}
Сначала понять предметную область,
а затем выбирать техническую реализацию.
\end{alertblock}

\end{frame}


% --- ER-модель ---
\begin{frame}{ER-модель}

\begin{block}{Определение}
\textbf{ER-модель} (Entity--Relationship Model) была
предложена Питером Ченом (Peter Chen) как концептуальная
модель данных, предназначенная для представления
семантики объектов предметной области и отношений
между ними при проектировании баз данных.
\end{block}

\vspace{0.4cm}

\begin{center}
\Large
Сущности
\quad+\quad
Атрибуты
\quad+\quad
Связи
\end{center}

Графическое представление — \textbf{ER-диаграмма (ERD)}.

\end{frame}


% --- Элементы ER ---
\begin{frame}{Базовые элементы ER-диаграммы (нотация Чена)}

\begin{table}
\centering
\begin{tabular}{>{\bfseries}p{0.3\textwidth}|p{0.45\textwidth}}
\toprule
Элемент & Значение\\
\midrule
Прямоугольник & сущность\\
Овал & атрибут\\
Ромб & связь\\
Линия & соединение элементов\\
Подчёркивание атрибута & первичный ключ\\
\bottomrule
\end{tabular}
\end{table}

\vspace{0.4cm}

ER-диаграмма позволяет проще обсуждать структуру
с заказчиком и находить ошибки \textbf{до} того,
как система начала работать.

\end{frame}


% --- Пример ER-модели (TikZ, Чен) ---
\begin{frame}{Пример ER-диаграммы (нотация Чена)}

\begin{center}
\resizebox{0.95\textwidth}{!}{%
\begin{tikzpicture}[node distance=0.8cm and 1cm]

\node[entity] (group) {\large{Группа}};
\node[relationship, right=of group] (r1) {\large{обучается}};
\node[entity, right=of r1] (student) {\large{Студент}};
\node[relationship, right=of student] (r2) {\large{получает}};
\node[entity, right=of r2] (grade) {\large{Оценка}};

\draw[thick] (group) -- (r1) node[card, midway] {\large{1}};
\draw[thick] (r1) -- (student) node[card, midway] {\large{N}};
\draw[thick] (student) -- (r2) node[card, midway] {\large{1}};
\draw[thick] (r2) -- (grade) node[card, midway] {\large{N}};

\node[attribute, below=1cm of group, xshift=-1.0cm] (a1) {\large{номер}};
\node[attribute, below=1cm of group, xshift=1.2cm] (a2) {\large{курс}};
\draw (group) -- (a1);
\draw (group) -- (a2);

\node[attribute, below=1cm of student, xshift=-0.9cm] (a3) {\large{\underline{номер}}};
\node[attribute, below=1cm of student, xshift=1.2cm] (a4) {\large{ФИО}};
\draw (student) -- (a3);
\draw (student) -- (a4);

\node[attribute, below=1cm of grade, xshift=-0.9cm] (a5) {\large{балл}};
\node[attribute, below=1cm of grade, xshift=1.2cm] (a6) {\large{дата}};
\draw (grade) -- (a5);
\draw (grade) -- (a6);

\end{tikzpicture}
}
\end{center}

\vspace{0.2cm}
\begin{itemize}
  \item прямоугольник — сущность, ромб — связь, овал — атрибут;
  \item подчёркивание атрибута — первичный ключ.
\end{itemize}

\end{frame}


% --- Crow's Foot ---
\begin{frame}{Нотация Crow's Foot}

Более компактная нотация,
часто используемая в практическом проектировании.

\begin{center}
\resizebox{0.72\textwidth}{!}{%
\begin{tikzpicture}
\node[entity, minimum width=3cm] (group) {Группа};
\node[entity, minimum width=3cm, right=6cm of group] (student) {Студент};

\draw[very thick, svfuDark] (group.east) -- (student.west);

\draw[very thick, svfuDark] ($(group.east)+(0.5,0.3)$) -- ($(group.east)+(0.5,-0.3)$);
\draw[very thick, svfuDark] ($(group.east)+(0.8,0.3)$) -- ($(group.east)+(0.8,-0.3)$);

\draw[very thick, svfuDark] ($(student.west)+(-0.7,0)$) -- ($(student.west)+(0,0.3)$);
\draw[very thick, svfuDark] ($(student.west)+(-0.7,0)$) -- ($(student.west)+(0,-0.3)$);
\draw[very thick, svfuDark] ($(student.west)+(-0.7,0)$) -- (student.west);

\draw[very thick, svfuDark, fill=white] ($(student.west)+(-1.05,0)$) circle (0.16);

\node[above, font=\small] at ($(group.east)+(0.65,0.45)$) {один};
\node[above, font=\small] at ($(student.west)+(-0.35,0.45)$) {ноль или много};
\end{tikzpicture}
}
\end{center}

\begin{itemize}
  \item двойная чёрточка — «ровно один»;
  \item окружность — необязательность связи («ноль»);
  \item «воронья лапка» — «много».
\end{itemize}

\end{frame}


% --- Типы связей ---
\begin{frame}{Типы связей (кардинальность)}

\begin{table}
\centering
\begin{tabular}{>{\bfseries}p{0.25\textwidth}|p{0.45\textwidth}}
\toprule
Тип & Пример\\
\midrule
1:1 & Студент — Пропуск (паспорт — человек)\\
1:N & Группа — Студенты\\
N:M & Студенты — Дисциплины\\
\bottomrule
\end{tabular}
\end{table}

\vspace{0.4cm}

На практике особенно важна связь
\textbf{многие-ко-многим} — она преобразуется
в отдельную (связующую) таблицу.

\end{frame}


% --- Полная ER-диаграмма курса ---
\begin{frame}{Полная ER-диаграмма: предметная область «Университет»}

\begin{center}
\resizebox{0.92\textwidth}{!}{%
\begin{tikzpicture}[node distance=2.6cm and 4.0cm]

\node[entity] (group) {Группа};
\node[entity, right=of group] (student) {Студент};
\node[entity, right=of student] (grade) {Оценка};
\node[entity, below=of grade] (discipline) {Дисциплина};
\node[entity, below=of group] (teacher) {Преподаватель};

\node[relationship] (r1) at ($(group)!0.5!(student)$) {\large{обучается}};
\node[relationship] (r2) at ($(student)!0.5!(grade)$) {\large{получает}};
\node[relationship] (r3) at ($(grade)!0.5!(discipline)$) {\large{по}};
\node[relationship] (r4) at ($(teacher)!0.5!(discipline)$) {\large{ведёт}};

\draw[thick] (group) -- (r1) node[card, midway]  {\large{1}};
\draw[thick] (r1) -- (student) node[card, midway]{\large{N}};
\draw[thick] (student) -- (r2) node[card, midway]{\large{1}};
\draw[thick] (r2) -- (grade) node[card, midway]  {\large{N}};
\draw[thick] (grade) -- (r3) node[card, midway]  {\large{N}};
\draw[thick] (r3) -- (discipline) node[card, midway]{\large{1}};
\draw[thick] (teacher) -- (r4) node[card, midway] {\large{1}};
\draw[thick] (r4) -- (discipline) node[card, midway] {\large{N}};

\end{tikzpicture}
}
\end{center}

\begin{center}
Эта диаграмма — результат анализа предметной области
и отправная точка для логического проектирования.
\end{center}

\end{frame}


% --- ER-модель vs модель хранения ---
\begin{frame}{ER-модель и модель хранения — не одно и то же}

\begin{center}
\begin{tikzpicture}[node distance=1cm and 1.6cm]
\node[entity] (dom) {Предметная\\область};
\node[entity, below=1.1cm of dom] (er) {ER-модель};
\node[catbox, below=1.3cm of er, xshift=-3.6cm] (t) {Таблицы};
\node[catbox, below=1.3cm of er] (j) {JSON};
\node[catbox, below=1.3cm of er, xshift=3.6cm] (f) {Файлы};
\draw[thick] (dom) -- (er);
\draw[thick] (er) -- (t); \draw[thick] (er) -- (j); \draw[thick] (er) -- (f);
\end{tikzpicture}
\end{center}

\begin{alertblock}{Важно}
ER-диаграмма не является таблицей и не определяет
автоматически использование SQL или NoSQL —
она описывает структуру и семантику предметной области.
\end{alertblock}

\end{frame}


% =========================================================
% ЧАСТЬ VII. ОТ МОДЕЛИ К РЕАЛИЗАЦИИ
% =========================================================

% --- ANSI/SPARC (TikZ) ---
\begin{frame}{Трёхуровневая архитектура ANSI/SPARC}

Это не этапы проектирования,
а архитектурная модель самой СУБД.

\begin{center}
\begin{tikzpicture}[node distance=0.6cm]
\node[layer] (ext) {Внешний уровень\\ \footnotesize представления пользователей};
\node[layer, below=of ext] (con) {Концептуальный уровень\\ \footnotesize логическая структура БД};
\node[layer, below=of con] (int) {Внутренний уровень\\ \footnotesize физическое хранение};
\draw[-{Latex[length=2.5mm]}, thick, svfuAccent] (ext) -- (con);
\draw[-{Latex[length=2.5mm]}, thick, svfuAccent] (con) -- (int);
\end{tikzpicture}
\end{center}

\begin{center}
Цель — \textbf{независимость данных}.
\end{center}

\end{frame}


% --- Три уровня подробнее ---
\begin{frame}{Три уровня подробнее}

\begin{itemize}
  \item \textbf{Внешний} — как данные видит конкретный
        пользователь или приложение (студент видит свои
        оценки, деканат — статистику факультета).
        В PostgreSQL этому соответствует \texttt{VIEW}.
  \item \textbf{Концептуальный} — единое логическое описание
        всей БД: сущности, атрибуты, связи, ограничения.
  \item \textbf{Внутренний} — как данные физически хранятся:
        файлы, страницы, индексы, буферы.
\end{itemize}

\vspace{0.3cm}

\begin{alertblock}{Главная идея}
Пользователю не нужно знать,
в каких физических страницах лежит строка таблицы.
\end{alertblock}

\end{frame}


% --- Независимость данных ---
\begin{frame}{Независимость данных}

\begin{table}
\centering
\begin{tabular}{p{0.3\textwidth}|p{0.58\textwidth}}
\toprule
\textbf{Тип} & \textbf{Что позволяет}\\
\midrule
Физическая &
изменять способы хранения и индексы
без переписывания приложений\\
Логическая &
изменять логическую структуру,
сохраняя существующие представления\\
\bottomrule
\end{tabular}
\end{table}

\end{frame}


% --- ER -> Table ---
\begin{frame}[fragile]{От ER-диаграммы к таблицам}

\begin{center}
\resizebox{0.85\textwidth}{!}{%
\begin{tikzpicture}
\node[entity] (g) {Группа};
\node[entity, right=2.2cm of g] (s) {Студент};
\draw[thick] (g) -- (s) node[card, midway, above] {\large{1 : N}};

\node[font=\Large, right=1.4cm of s] (arrow) {$\Rightarrow$};

\node[entity, fill=white, right=1.4cm of arrow, yshift=0.7cm] (t1) {student\_group};
\node[entity, fill=white, below=0.6cm of t1] (t2) {student};
\draw[-{Latex[length=2mm]}, thick, svfuAccent] (t2.north) -- (t1.south)
  node[midway, right, font=\tiny] {\large{group\_id (KEY)}};
\end{tikzpicture}
}
\end{center}

\small
\begin{verbatim}
CREATE TABLE student_group (
    id SERIAL PRIMARY KEY,
    name TEXT
);

CREATE TABLE student (
    id SERIAL PRIMARY KEY,
    name TEXT,
    group_id INTEGER REFERENCES student_group(id)
);
\end{verbatim}

\end{frame}


% --- Ограничения целостности ---
\begin{frame}{Ограничения целостности}

Ограничения не позволяют БД
перейти в некорректное состояние.

\vspace{0.4cm}

\begin{itemize}
  \item \texttt{PRIMARY KEY} — уникальность и обязательность;
  \item \texttt{FOREIGN KEY} — ссылочная целостность;
  \item \texttt{NOT NULL} — обязательность значения;
  \item \texttt{UNIQUE} — уникальность значения;
  \item \texttt{CHECK} — произвольное условие на значение.
\end{itemize}

\vspace{0.3cm}

\begin{alertblock}{Идея}
Правила должны контролироваться самой БД,
а не только программой пользователя.
\end{alertblock}

\end{frame}


% --- Слайд 1: Проблема ---
\begin{frame}{Нормализация (введение)}

\textbf{Нормализация} — процесс организации данных,
направленный на уменьшение избыточности
и устранение аномалий.

\vspace{0.1cm}

\begin{alertblock}{Плохой вариант}
\begin{center}
\begin{tabular}{|l|l|l|l|}
\hline
\textbf{student} & \textbf{name} & \textbf{discipline} & \textbf{teacher} \\
\hline
1 & Иванов & Математика & Сидоров \\
1 & Иванов & Программирование & Козлов \\
2 & Петров & Математика & Сидоров \\
2 & Петров & Физика & Орлова \\
\hline
\end{tabular}
\end{center}
\end{alertblock}

\vspace{-0.3cm}

\begin{block}{Проблема}
Один и тот же факт (преподаватель, название дисциплины) повторяется много раз.\\[0.2cm]
Изменение одного факта требует изменения множества строк → риск противоречий и аномалий.
\end{block}

\end{frame}



% --- Слайд 2: Решение ---
\begin{frame}{Нормализация}
\vspace{-0.5cm} 
\begin{columns}[T]
\begin{column}{0.48\textwidth}
\begin{block}{Нормализованные таблицы}
\small
\begin{tabular}{|l|l|}
\hline
\textbf{student} & \\
\hline
id & name \\
1 & Иванов \\
2 & Петров \\
\hline
\end{tabular}

\vspace{0.1cm}

\begin{tabular}{|l|l|}
\hline
\textbf{discipline} & \\
\hline
id & title \\
1 & Математика \\
2 & Программирование \\
3 & Физика \\
\hline
\end{tabular}

\vspace{0.1cm}

\begin{tabular}{|l|l|l|}
\hline
\textbf{grade} & & \\
\hline
student\_id & disc\_id & score \\
1 & 1 & 5 \\
1 & 2 & 5 \\
2 & 1 & 4 \\
2 & 3 & 3 \\
\hline
\end{tabular}
\end{block}
\end{column}

\begin{column}{0.48\textwidth}
\begin{exampleblock}{Преимущества}
\begin{itemize}
  \item Нет избыточности
  \item Изменение одного факта — в одном месте
  \item Легче поддерживать целостность
  \item Уменьшается риск аномалий
\end{itemize}
\end{exampleblock}

\vspace{0.4cm}

\begin{alertblock}{Важно}
Выделяют следующие нормальные формы: 1НФ,2НФ, 3НФ и т.д.
\end{alertblock}
\end{column}
\end{columns}

\end{frame}



% --- Access vs PostgreSQL ---
\begin{frame}{SQL, Access и PostgreSQL}

SQL объединяет различные средства работы
с реляционными БД (DDL, DML и другие).
SQL не заменяет проектирование.

\vspace{0.3cm}
\small
\begin{table}
\centering
\begin{tabular}{p{0.28\textwidth}|p{0.28\textwidth}|p{0.28\textwidth}}
\toprule
 & \textbf{Access} & \textbf{PostgreSQL}\\
\midrule
Тип & настольная СУБД & серверная СУБД\\
Интерфейс & сильный GUI & SQL + инструменты\\
Масштаб & небольшие системы & крупные системы\\
\bottomrule
\end{tabular}
\end{table}

\normalsize

\end{frame}


% --- Инструменты ---
\begin{frame}{Средства проектирования и администрирования}

\begin{itemize}
  \item \textbf{draw.io / diagrams.net} — универсальный
        инструмент построения ER-диаграмм;
  \item \textbf{pgModeler} — моделирование PostgreSQL;
  \item \textbf{DBeaver} — работа с различными СУБД
        и ER-диаграммами;
  \item \textbf{pgAdmin} — графическое администрирование PostgreSQL;
  \item \textbf{psql} — консольный клиент PostgreSQL (CLI).
\end{itemize}

\end{frame}


% =========================================================
% ЧАСТЬ VIII. ИТОГИ
% =========================================================

% --- Итоговая схема ---

\begin{frame}{Итоговая схема проектирования}


\vspace{0.6cm}

\begin{center}

\begin{tikzpicture}[node distance=1.2cm and 0.8cm]

\node[phase] (domain) {Предметная\\область};
\node[phase, right=of domain] (model) {ER-модель};
\node[phase, right=of model] (design) {Логическая\\модель};
\node[phase, right=of design] (db) {База данных};

\draw[-{Latex[length=2.5mm]}, thick, svfuAccent]
    (domain) -- (model);

\draw[-{Latex[length=2.5mm]}, thick, svfuAccent]
    (model) -- (design);

\draw[-{Latex[length=2.5mm]}, thick, svfuAccent]
    (design) -- (db);

\end{tikzpicture}

\end{center}

\vspace{0.5cm}

\begin{alertblock}{Главная идея}

Сначала мы понимаем что нужно хранить, затем строим модель и только после этого реализуем базу данных.

\end{alertblock}

\end{frame}


% --- Аналитический вывод ---
\begin{frame}{Аналитический вывод: от данных к модели хранения}

\begin{itemize}
  \item Выбор способа хранения — \textbf{результат} анализа
        данных, а не исходная точка проектирования.
  \item Табличное, документное и файловое представления
        не исключают друг друга и могут сочетаться
        в одной системе.
  \item ER-модель — семантическое описание предметной
        области, не привязанное к конкретной технологии.
  \item Разделение уровней (внешний / концептуальный /
        внутренний) обеспечивает независимость данных
        и гибкость системы при её росте.
\end{itemize}

\end{frame}


% % --- Задание ---
% \begin{frame}{Задание к следующей лекции}

% Выберите предметную область.

% \vspace{0.3cm}

% Например:
% \begin{itemize}
%   \item университет;
%   \item магазин;
%   \item библиотека;
%   \item кинотеатр.
% \end{itemize}

% \vspace{0.3cm}

% Определите:
% \begin{enumerate}
%   \item сущности и их атрибуты;
%   \item связи между сущностями и их кардинальность;
%   \item какие сведения структурированы, а какие —
%         естественнее представить документом или файлом.
% \end{enumerate}

% \vspace{0.3cm}

% \textbf{Результат:} черновая ER-диаграмма.

% \end{frame}





% --- Спасибо ---
\begin{frame}[plain,noframenumbering]

\thispagestyle{empty}

\vspace*{\fill}

\begin{center}

{\LARGE\bfseries
Спасибо за внимание!
}

\end{center}

\vspace*{\fill}

\end{frame}


\end{document}
